\medskip

%Une urne contient 21 jetons numérotés de 1 
%à 21 indiscernables au toucher.
%
%On tire un jeton au hasard.

\begin{enumerate}
\item %On note $A$ l’évènement \og obtenir 2, 3 ou 10 \fg.
Il y a trois chances de succès sur 21 sorties différentes, donc $p(A) = \dfrac{3}{21} = \dfrac{3 \times 1}{3 \times 7} = \dfrac17$.

%Calculer la probabilité de l’évènement $A$.

%On donnera le résultat sous forme de fraction irréductible.
\item 
	\begin{enumerate}
		\item %On note $B$ l’évènement \og obtenir un jeton dont le numéro est un diviseur de 24 \fg. 

%Donner les issues de l’évènement $B$.
On a $24 = 1\times 24 = 2 \times 12 = 3 \times 8 = 4 \times 6$.

Les diviseurs de 24 sont donc dans l'ordre croissant : 1~;~2~;~3~;~4~;~6~;~8~;~12~;~24.

Les jetons de l'unité qui portent un nombre diviseur de 24 sont donc :

1~;~2~;~3~;~4~;~6~;~8~;~12
		\item %Déterminer la probabilité de l’évènement $B$.
		Il y a ici 7 possibilités de gagner, donc $p(B) = \dfrac{7}{21} = \dfrac{7 \times 1}{7 \times 3} = \dfrac13$.
	\end{enumerate}
\end{enumerate}

\medskip
